\documentclass[11pt, a4paper]{article}

% Gemeinsame Style-Definitionen (TEXINPUTS muss templates/ enthalten — siehe scripts/build_tex.py)
% bfw-style.tex — gemeinsame Style-Definitionen für alle Handouts/Aufgabenhefte
% Voraussetzung: Kompilation mit xelatex (wegen fontspec)

\usepackage[ngerman]{babel}
\usepackage{geometry}
\geometry{a4paper, top=2.2cm, bottom=2.2cm, left=2.3cm, right=2.3cm}

\usepackage{fontspec}
% Fallback-Kette: Source Sans Pro -> Calibri -> Segoe UI -> Latin Modern Sans
% (Latin Modern ist immer mit MiKTeX vorhanden)
\IfFontExistsTF{Source Sans Pro}{%
  \setmainfont{Source Sans Pro}[Ligatures=TeX]%
}{\IfFontExistsTF{Calibri}{%
  \setmainfont{Calibri}[Ligatures=TeX]%
}{\IfFontExistsTF{Segoe UI}{%
  \setmainfont{Segoe UI}[Ligatures=TeX]%
}{%
  \setmainfont{Latin Modern Sans}[Ligatures=TeX]%
}}}
\IfFontExistsTF{Source Code Pro}{%
  \setmonofont{Source Code Pro}[Scale=0.92]%
}{\IfFontExistsTF{Consolas}{%
  \setmonofont{Consolas}[Scale=0.92]%
}{%
  \setmonofont{Latin Modern Mono}[Scale=0.92]%
}}

\usepackage{xcolor}
% Farbpalette: hell, freundlich, modern
\definecolor{bfwPrimary}{HTML}{0EA5A4}    % Petrol/Türkis - Primär
\definecolor{bfwPrimaryDark}{HTML}{0F766E} % dunkler Petrol für Headings
\definecolor{bfwAccent}{HTML}{F59E0B}     % Amber - Akzent
\definecolor{bfwSuccess}{HTML}{10B981}    % Grün - Erfolg / Lösung
\definecolor{bfwWarn}{HTML}{F97316}       % Orange - Achtung
\definecolor{bfwInk}{HTML}{0F172A}        % fast Schwarz
\definecolor{bfwInkSoft}{HTML}{475569}    % gedämpftes Grau
\definecolor{bfwBgSoft}{HTML}{F8FAFC}     % off-white
\definecolor{bfwBgTeal}{HTML}{ECFEFF}     % helles Türkis
\definecolor{bfwBgAmber}{HTML}{FEF3C7}    % helles Gelb
\definecolor{bfwBgRose}{HTML}{FFE4E6}     % helles Rosé für Eselsbrücken

\usepackage[colorlinks=true, linkcolor=bfwPrimaryDark, urlcolor=bfwPrimaryDark]{hyperref}
\usepackage{enumitem}
\setlist[itemize]{leftmargin=*, itemsep=2pt, topsep=2pt}
\setlist[enumerate]{leftmargin=*, itemsep=2pt, topsep=2pt}

\usepackage[most]{tcolorbox}
\usepackage{booktabs}
\usepackage{tikz}
\usetikzlibrary{decorations.pathmorphing, positioning, shapes.geometric, arrows.meta}

% Merkkästchen — wichtige Definition / Kernpunkt
\newtcolorbox{merksatz}[1][]{%
  enhanced, breakable,
  colback=bfwBgTeal,
  colframe=bfwPrimary,
  boxrule=0.6pt,
  arc=4pt,
  left=10pt, right=10pt, top=8pt, bottom=8pt,
  fonttitle=\bfseries\color{bfwPrimaryDark},
  title={\faLightbulb\ Merksatz},
  #1
}

% Eselsbrücke — Mnemoniks
\newtcolorbox{eselsbruecke}[1][]{%
  enhanced, breakable,
  colback=bfwBgRose,
  colframe=bfwAccent,
  boxrule=0.6pt,
  arc=4pt,
  left=10pt, right=10pt, top=8pt, bottom=8pt,
  fonttitle=\bfseries\color{bfwWarn},
  title={Eselsbrücke},
  #1
}

% Beispiel-Box
\newtcolorbox{beispiel}[1][]{%
  enhanced, breakable,
  colback=bfwBgSoft,
  colframe=bfwInkSoft,
  boxrule=0.4pt,
  arc=3pt,
  left=10pt, right=10pt, top=8pt, bottom=8pt,
  fonttitle=\bfseries\color{bfwInk},
  title={Beispiel},
  #1
}

% Lösungs-Box (für Aufgabenhefte)
\newtcolorbox{loesung}[1][]{%
  enhanced, breakable,
  colback=bfwBgAmber!40,
  colframe=bfwSuccess,
  boxrule=0.6pt,
  arc=3pt,
  left=10pt, right=10pt, top=8pt, bottom=8pt,
  fonttitle=\bfseries\color{bfwSuccess!70!black},
  title={Lösung},
  #1
}

% Glossar-Eintrag
\newcommand{\glossareintrag}[2]{%
  \par\noindent\textbf{\color{bfwPrimaryDark}#1} \enspace #2\par\smallskip
}

% Section-Stil — etwas farbiger als Standard
\usepackage{titlesec}
\titleformat{\section}
  {\Large\bfseries\color{bfwPrimaryDark}}
  {\thesection}{0.6em}{}
\titleformat{\subsection}
  {\large\bfseries\color{bfwPrimary}}
  {\thesubsection}{0.5em}{}
\titleformat{\subsubsection}
  {\normalsize\bfseries\color{bfwInk}}
  {\thesubsubsection}{0.4em}{}

% TikZ-Stil "skizzenhaft" — etwas weicher als Rough.js, aber stabil
\tikzset{
  roughbox/.style = {
    draw=bfwPrimaryDark, thick, fill=bfwBgTeal,
    rounded corners=4pt, inner sep=6pt
  },
  roughaccent/.style = {
    draw=bfwAccent, thick, fill=bfwBgAmber,
    rounded corners=4pt, inner sep=6pt
  },
  roughline/.style = {
    draw=bfwInkSoft, thick, -{Stealth[length=2.5mm]}
  }
}

% Bullet-Symbole (bewusst kein FontAwesome — vermeidet Font-Installations-Probleme)
\newcommand{\faLightbulb}{$\bullet$}
\newcommand{\faBookmark}{$\bullet$}

% Header/Footer
\usepackage{fancyhdr}
\pagestyle{fancy}
\fancyhf{}
\renewcommand{\headrulewidth}{0pt}
\fancyfoot[L]{\small\color{bfwInkSoft}\thepage}
\fancyfoot[R]{\small\color{bfwInkSoft} BFW M-064 Beschaffung im Büromanagement}


\title{\vspace*{-1.5cm}\Huge\bfseries\color{bfwPrimaryDark}Tag~1: Büroraumplanung \& Büroformen\\[0.4em]
  \large\color{bfwInkSoft}\textnormal{Wie der Raum die Arbeit prägt}}
\author{\textsf{Büroprozesse gestalten und Arbeitsvorgänge organisieren \textperiodcentered{} M-067}\\
\textsf{\small\copyright\ Can Siebert\,\textperiodcentered\,2026}}
\date{29.06.2026}

\begin{document}
\maketitle
\thispagestyle{empty}

\vspace{1em}
\begin{merksatz}[title={\bfseries Worum geht es heute?}]
Das Büro ist mehr als ein Raum mit Schreibtischen --- es ist das Werkzeug, mit dem
Büroarbeit gelingt oder scheitert. Eine durchdachte Raumplanung steigert Produktivität,
Gesundheit und Zufriedenheit; eine schlechte verursacht Lärm, Krankheit und Kosten.
In diesem Handout lernen Sie, wie Büroflächen geplant werden, aus welchen Teilflächen
sich ein Arbeitsplatz zusammensetzt und welche Büroformen es gibt --- von der klassischen
Zelle bis zum modernen Desksharing. Das ist die Grundlage für den gesamten
Modulverlauf und für die IHK-Abschlussprüfung.
\end{merksatz}

\tableofcontents
\vspace{1em}
\hrule
\vspace{1em}

\section{Warum Büroraumplanung wichtig ist}

Die Raumgestaltung eines Büros wirkt sich entscheidend auf die \emph{Produktivität}
und die \emph{Arbeitszufriedenheit} seiner Nutzer aus. Deshalb sollten die vorhandenen
räumlichen Ressourcen eines Unternehmens sorgfältig geplant und gestaltet werden ---
schließlich verursachen Büroflächen erhebliche Kosten (Miete, Energie, Reinigung) und
binden Kapital. Ein Büro ist dabei definiert als ein Raum, der überwiegend zur
Erledigung von Verwaltungstätigkeiten dient.

Als Leitfaden dient die \emph{DGUV Information 215-441 „Büroraumplanung, Hilfen für das
systematische Planen und Gestalten von Büros”}. Sie konkretisiert die sicherheits\-tech\-nischen,
arbeitsmedizinischen, ergonomischen und arbeitspsychologischen Anforderungen, die bei
der Planung von Flächen von Büroräumen und Bürobereichen zu beachten sind. Wer ein Büro
plant, plant also nicht nur Möbel, sondern auch Gesundheit, Sicherheit und
Konzentrationsfähigkeit der Beschäftigten.

\begin{merksatz}
Gute Büroraumplanung ist kein Luxus, sondern Wirtschaftlichkeit: Zufriedene, gesunde
Beschäftigte arbeiten produktiver, sind seltener krank und bleiben dem Unternehmen
länger erhalten.
\end{merksatz}

\section{Die sechs Teilflächen eines Büros}

Die Fläche eines Büros lässt sich in sechs verschiedene \emph{Teilflächen} untergliedern.
Erst die Summe dieser Teilflächen ergibt den tatsächlichen Platzbedarf eines
Arbeitsplatzes --- man kann den Bedarf also nicht einfach schätzen, sondern muss ihn
aus seinen Bestandteilen aufbauen.

\begin{enumerate}
  \item \textbf{Stellfläche} --- die benötigte Bodenfläche für die Einrichtung, z.\,B.\ für
    Schreibtisch, Schrank, Regal und Pflanzen.
  \item \textbf{Bewegungsfläche} --- freie, unverstellte Bodenfläche am Arbeitsplatz, die
    wechselnde Arbeitshaltungen und Ausgleichsbewegungen ermöglicht.
  \item \textbf{Funktionsfläche} --- die Fläche, die z.\,B.\ beim Öffnen von Schranktüren
    oder Schubladen überdeckt wird.
  \item \textbf{Benutzerfläche am Arbeitsplatz} --- die Fläche am persönlich zugewiesenen
    Arbeitsplatz, die der Nutzer bei seiner jeweiligen Tätigkeit benötigt.
  \item \textbf{Verkehrswegefläche} --- die Fläche, die innerhalb eines Büros für den
    Personenverkehr und den Transport von Material notwendig ist.
  \item \textbf{Besucher- und Besprechungsfläche} --- Benutzerfläche an Besucher- und
    Besprechungsplätzen, Schränken usw.\ im Bürobereich.
\end{enumerate}

\begin{eselsbruecke}
\textbf{„Steh Bequem, Funktioniere, Benutze Verkehr Besuchend”} --- die Anfangs\-buch\-staben
\emph{St-Be-Fu-Be-Ve-Be} erinnern an \emph{Stell-, Bewegungs-, Funktions-, Benutzer-,
Verkehrswege- und Besucherfläche}. Wichtig ist vor allem das Prinzip: Der Platzbedarf
entsteht aus der \emph{Addition} aller sechs Flächen.
\end{eselsbruecke}

\subsection{Quadratmeterbedarf je Arbeitsplatz}

Aus der Addition der jeweils ermittelten Werte ergibt sich der \emph{Quadratmeterbedarf
pro Arbeitsplatz}. Als Richtwerte gelten:

\begin{center}
\begin{tabular}{p{7cm} r}
\toprule
\textbf{Arbeitsplatztyp} & \textbf{Flächenbedarf}\\
\midrule
Bildschirmarbeitsplatz (Einzel) & 8--10 m\textsuperscript{2}\\
Arbeitsplatz im Großraumbüro & ca.\ 12--15 m\textsuperscript{2}\\
\bottomrule
\end{tabular}
\end{center}

\medskip
Warum braucht ein Arbeitsplatz im Großraumbüro \emph{mehr} Fläche als ein einzelner
Bildschirmarbeitsplatz? Weil dort mehr Verkehrswege notwendig sind: Je mehr Menschen
sich einen Raum teilen, desto mehr Fläche entfällt anteilig auf die Wege zwischen den
Arbeitsplätzen.

\begin{beispiel}
Ein Raum von 60~m\textsuperscript{2} soll mit einzelnen Bildschirmarbeitsplätzen zu je
10~m\textsuperscript{2} ausgestattet werden. Rechnerisch passen
$60 \div 10 = 6$ Arbeitsplätze hinein. Würde man denselben Raum als Großraum\-bereich
mit 12~m\textsuperscript{2} je Platz planen, wären es nur fünf Plätze --- die zusätzlichen
Verkehrswege kosten Fläche.
\end{beispiel}

\section{Büroformen im Überblick}

Büroflächen können sehr unterschiedlich gestaltet sein, deshalb unterscheidet man
verschiedene \emph{Büroformen}. Keine davon ist universell „die beste” --- jede hat
Stärken und Schwächen, die zur jeweiligen Tätigkeit und Unternehmenskultur passen
müssen.

\subsection{Das Zellenbüro}

Das \textbf{Zellenbüro} (1--3 Personen) ist die klassische Büroform und in deutschen
Unternehmen oft anzutreffen. Es erfreut sich großer Beliebtheit, weil die Mitarbeitenden
ihr Büro individuell gestalten und Umgebungseinflüsse wie Beleuchtung und Raumklima
selbst regulieren können. Im Einzelbüro lässt sich konzentriert arbeiten und es können
vertrauliche Gespräche geführt werden. Die Kommunikation mit anderen erfolgt über
technische Hilfsmittel (Telefon, E-Mail) oder in Besprechungen. Nachteil: Der Nutzer
fühlt sich eventuell isoliert, und der Flächenbedarf ist hoch. Einzelbüros haben heute
vor allem Führungskräfte; in Unternehmen mit überwiegend Großraumbüros gelten sie als
\emph{Statussymbol}.

\subsection{Das Gruppenbüro}

\textbf{Gruppenbüros} beherbergen drei bis ca.\ 25 Arbeitsplätze und werden meist von
einer Abteilung oder einem Team gemeinsam genutzt --- ideal für projektorientiertes
Arbeiten. Sie bestehen nicht aus eigenständigen Räumen, sondern aus einem großen Raum,
der durch Trennwände oder Raumgliederungselemente unterteilt wird.

Diese Büroform fördert eine \emph{intensive Kommunikation und Kooperation} und ermöglicht
einen schnellen Informationsfluss bei kurzen Wegen. Durch die ständige Kommunikation
entsteht hohe Transparenz, was die Vertretung von Kollegen im Urlaubs- oder
Krankheitsfall erleichtert. Nachteil: Bei vielen Personen entsteht ein hoher
\emph{Geräuschpegel}, der konzentriertes Arbeiten erschwert. Abhilfe schaffen
\emph{Telefon- und Arbeitszellen} als Rückzugsmöglichkeit sowie schallabsorbierende
Decken und schalldämpfende Teppiche. Der Übergang zum Großraumbüro ist fließend.

\subsection{Das Großraumbüro}

Ein \textbf{Großraumbüro} nimmt oft eine ganze Etage ein (Flächen von 400~m\textsuperscript{2}
bis zu mehreren tausend m\textsuperscript{2}) und bietet Arbeitsplätze für 25 bis mehrere
hundert Mitarbeitende. Der Raum wird meist durch mittelhohe Stellwände, Schränke,
Wandsysteme und Pflanzen gegliedert und bietet hohe \emph{Flexibilität}: Veränderungen
der Mitarbeiterzahl oder Teamzugehörigkeit lassen sich ohne großen Aufwand umsetzen.

Für diese Büroform ist eine \emph{Klimaanlage} erforderlich, wodurch der einzelne
Mitarbeiter das Klima nicht selbst beeinflussen kann. Die mangelnde Privatsphäre und
das Gefühl der ständigen Kontrolle sind zwei große Probleme. Gefördert wird hingegen
der \emph{informelle Austausch} von Wissen. Das Großraumbüro gilt als besonders
\emph{kostengünstig}, da die Bürofläche effizient genutzt wird und teure Geräte (z.\,B.\
Drucker) in geringerer Zahl angeschafft werden müssen. Diese Ersparnis wird laut Studien
jedoch häufig durch Unzufriedenheit, höheren Krankenstand und geringere Produktivität
wieder aufgehoben. Viele Unternehmen nutzen Großraumbüros vor allem für Call-Center.

\subsection{Das Kombibüro}

Das \textbf{Kombibüro} besteht aus kleinen Einzelbüros an der Fensterfront des Gebäudes
und hat anstatt eines Flurs in der Mitte eine sogenannte \emph{Multifunktionszone}, die
jederzeit veränderbar ist. Diese Zone besteht oft aus Besprechungs-, Archiv- und
Pausenzonen sowie temporären Arbeitsplätzen; auch ein Technikpool (Kopierer, Drucker)
findet hier Platz. Glaswände und -türen lassen natürliches Licht in die Multifunktions\-zone,
bewirken aber auch, dass sich Mitarbeitende beobachtet fühlen und die visuellen
Störungen steigen.

Das Kombibüro \emph{vereint die Vorteile} von Zellen- und Großraumbüro und eignet sich
vor allem für Arbeitsplätze, bei denen störungsfreie Einzelarbeit und gelegentliche
Teamarbeit erforderlich sind. Die Einzelbüros sind durch raumhohe, versetzbare Wände
getrennt, was einen schnellen Umbau von Einzel- zu Doppelbüros ermöglicht. Wichtig:
Die Einzelbüros dürfen nicht zu klein zugeschnitten werden.

\subsection{Vergleich der Büroformen}

\begin{tabular}{p{2.6cm} p{4.7cm} p{4.7cm}}
\toprule
\textbf{Büroform} & \textbf{Vorteile} & \textbf{Nachteile}\\
\midrule
Zellenbüro & Konzentration, Vertraulichkeit, individuelle Gestaltung & Isolationsgefühl, hoher Flächenbedarf\\
Gruppenbüro & Kommunikation, schneller Informationsfluss, leichte Vertretung & Hoher Geräuschpegel, weniger Privatsphäre\\
Großraumbüro & Kostengünstig, flexibel, informeller Wissensaustausch & Mangelnde Privatsphäre, Kontrollgefühl, höherer Krankenstand\\
Kombibüro & Einzelarbeit \emph{und} Teamarbeit, flexibel umbaubar & Beobachtungsgefühl durch Glas, höhere Baukosten\\
\bottomrule
\end{tabular}

\medskip
\begin{eselsbruecke}
\textbf{„Zelle = Ruhe, Gruppe = Reden, Großraum = Sparen, Kombi = Beides.”}
So merken Sie sich den jeweiligen Schwerpunkt der vier klassischen Büroformen.
\end{eselsbruecke}

\section{Moderne und flexible Bürokonzepte}

Die Arbeitswelt verändert sich: mehr Projektarbeit, mehr Austausch, mehr Mobilität.
Daraus sind Bürokonzepte entstanden, die über die vier klassischen Formen hinausgehen.

\subsection{Das reversible Büro}

Bei der Planung von Bürogebäuden wird immer stärker auf \emph{Variabilität und
Flexibilität} geachtet. Das \textbf{reversible Büro} kann mit geringem Aufwand an die
jeweiligen organisatorischen Erfordernisse angepasst werden. Ausgangspunkt ist ein
Bauwerk, in dem sowohl Ein- bis Mehrpersonenzellenbüros als auch Gruppen- oder
Kombibüros Platz finden. Feste Wände sind die Ausnahme; die Aufteilung erfolgt durch
flexible Wandsysteme. Möbel, elektrische Ausstattung und Beleuchtung müssen sich leicht
ändern lassen. So kann jede Etage anforderungsgerecht eine andere Büroform aufnehmen ---
und das Unternehmen kann schnell und kostengünstig auf Veränderungen reagieren.

\subsection{Das non-territoriale Büro (Desksharing)}

Leere Büros und ungenutzte Arbeitsplätze sind ein hoher Kostenfaktor: Viele Beschäftigte
nutzen ihren Arbeitsplatz nur selten, weil sie wegen Urlaub, Krankheit, Außendienst,
Projektarbeit, Weiterbildung oder Besprechungen abwesend sind. Mobile Technik macht das
\textbf{non-territoriale Büro} (Desksharing) möglich.

\begin{merksatz}
Beim non-territorialen Büro ist die Anzahl der Arbeitsplätze \emph{kleiner} als die Zahl
der Beschäftigten. Nicht jeder hat einen persönlichen Schreibtisch --- bei Arbeitsbeginn
sucht man sich den nächsten freien Platz, oft unterstützt durch eine Buchungssoftware.
\end{merksatz}

So wird eine bessere Auslastung der Fläche erreicht und es lassen sich erhebliche Kosten
einsparen (Büroflächen, Möbel, Technik, Reinigung, Energie). Fachliche Unterlagen werden
digital gespeichert oder zentral abgelegt; persönliche Unterlagen kommen in einen
fahrbaren Rollcontainer, der bei Arbeitsbeginn zum freien Schreibtisch gebracht wird.
Besonders geeignet ist das Konzept wegen ihrer Kommunikationsstrukturen für Gruppen-
und Kombibüros. Zu beachten ist: Die Arbeitsplätze müssen leicht an die Körpergröße des
Nutzers anpassbar sein, und Kunden, Besucher oder Kollegen anderer Abteilungen haben
keine feste Anlaufstelle mehr --- sie müssen erst prüfen, wo ihr Ansprechpartner sitzt.

\subsection{Open-Space, Homeoffice und New Work}

Verwandt mit dem Großraum- und non-territorialen Büro ist das \emph{Open-Space-Konzept},
das offene, kommunikationsfördernde Flächen mit unterschiedlichen Zonen kombiniert
(Fokuszonen, Begegnungszonen, Rückzugsräume). Treiber ist die Idee des \emph{New Work}:
selbstbestimmtes, vernetztes Arbeiten.

Gleichzeitig hat sich \emph{Homeoffice} bzw.\ \emph{mobiles Arbeiten} stark verbreitet.
Wenn Beschäftigte zeitweise von zu Hause oder unterwegs arbeiten, sind weniger
Arbeitsplätze gleichzeitig belegt. Das senkt den Flächenbedarf und macht Konzepte wie
Desksharing wirtschaftlich besonders attraktiv --- erfordert aber klare Regeln,
verlässliche Technik und eine durchdachte Raumplanung, damit Zusammenarbeit und
Unternehmenskultur nicht leiden.

\section{Zusammenfassung}

\begin{enumerate}
  \item Büroraumplanung beeinflusst \textbf{Produktivität, Gesundheit und Zufriedenheit}; Leitfaden ist die DGUV Information 215-441.
  \item Ein Arbeitsplatz besteht aus \textbf{sechs Teilflächen}: Stell-, Bewegungs-, Funktions-, Benutzer-, Verkehrswege- und Besucher-/Besprechungsfläche.
  \item Der \textbf{Flächenbedarf} ergibt sich aus der Addition dieser Teilflächen: Bildschirmarbeitsplatz 8--10 m\textsuperscript{2}, Großraum ca.\ 12--15 m\textsuperscript{2}.
  \item Vier klassische \textbf{Büroformen}: Zellen- (Ruhe), Gruppen- (Kommunikation), Großraum- (Kosten), Kombibüro (beides).
  \item Moderne Konzepte: \textbf{reversibles Büro} (flexibel umbaubar) und \textbf{non-territoriales Büro/Desksharing} (weniger Plätze als Beschäftigte).
  \item \textbf{Open-Space, Homeoffice und New Work} verändern den Flächenbedarf und fördern flexible Konzepte.
\end{enumerate}

\newpage

\section*{Glossar}
\addcontentsline{toc}{section}{Glossar}

\glossareintrag{Büroraumplanung}{Systematische Planung und Gestaltung von Büroflächen unter sicherheits\-technischen, ergonomischen und arbeitspsychologischen Gesichtspunkten.}

\glossareintrag{DGUV Information 215-441}{Leitfaden „Büroraumplanung”, der die Anforderungen an die Planung von Büroflächen konkretisiert.}

\glossareintrag{Teilflächen}{Die sechs Bestandteile, aus denen sich der Flächenbedarf eines Arbeitsplatzes zusammensetzt.}

\glossareintrag{Stellfläche}{Bodenfläche für die Einrichtung wie Schreibtisch, Schrank, Regal und Pflanzen.}

\glossareintrag{Bewegungsfläche}{Freie, unverstellte Fläche am Arbeitsplatz für wechselnde Arbeitshaltungen und Ausgleichsbewegungen.}

\glossareintrag{Funktionsfläche}{Fläche, die beim Öffnen von Schranktüren oder Schubladen überdeckt wird.}

\glossareintrag{Verkehrswegefläche}{Fläche für Personenverkehr und Materialtransport innerhalb des Büros.}

\glossareintrag{Zellenbüro}{Klassische Büroform für 1--3 Personen; konzentriertes und vertrauliches Arbeiten, hoher Flächenbedarf.}

\glossareintrag{Gruppenbüro}{Büroform für 3 bis ca.\ 25 Arbeitsplätze; fördert Kommunikation, aber hoher Geräuschpegel.}

\glossareintrag{Großraumbüro}{Büroform für 25 bis mehrere hundert Personen auf großer Fläche; kostengünstig, aber mangelnde Privatsphäre.}

\glossareintrag{Kombibüro}{Einzelbüros an der Fensterfront plus zentrale Multifunktionszone; vereint Einzel- und Teamarbeit.}

\glossareintrag{Multifunktionszone}{Veränderbare Zone in der Mitte eines Kombibüros mit Besprechungs-, Archiv- und Pausenbereichen.}

\glossareintrag{Reversibles Büro}{Büroform, die mit geringem Aufwand durch flexible Wandsysteme an neue Erfordernisse angepasst werden kann.}

\glossareintrag{Non-territoriales Büro}{Desksharing-Konzept; es gibt weniger Arbeitsplätze als Beschäftigte, kein persönlicher Schreibtisch.}

\glossareintrag{Desksharing}{Mehrere Beschäftigte teilen sich Arbeitsplätze, die bei Bedarf gebucht und neu eingestellt werden.}

\glossareintrag{Open-Space}{Offenes Bürokonzept mit verschiedenen Zonen für Fokus, Begegnung und Rückzug.}

\glossareintrag{Mobiles Arbeiten}{Arbeit außerhalb der Betriebsstätte (z.\,B.\ Homeoffice); senkt den gleichzeitigen Flächenbedarf.}

\glossareintrag{Statussymbol}{Hier: das Einzelbüro als Zeichen einer höheren Hierarchieebene in Unternehmen mit Großraumbüros.}

\end{document}
